\documentclass[11pt]{article}
\usepackage[margin=1in]{geometry}
\usepackage{amsmath,amssymb,amsthm}
\usepackage{booktabs}
\usepackage{array}
\usepackage{graphicx}
\usepackage{xcolor}
\usepackage[colorlinks=true,linkcolor=blue!60!black,citecolor=blue!60!black,urlcolor=blue!60!black]{hyperref}

\newtheorem{proposition}{Proposition}
\newtheorem{lemma}{Lemma}
\newtheorem{corollary}{Corollary}
\theoremstyle{remark}
\newtheorem{remark}{Remark}

\newcommand{\R}{\mathbb{R}}
\newcommand{\E}{\mathbb{E}}
\newcommand{\Var}{\operatorname{Var}}
\newcommand{\ind}[1]{\mathbf{1}\{#1\}}
\newcommand{\bx}{\mathbf{x}}
\newcommand{\Sig}{\Sigma}
\newcommand{\Tk}{\mathcal{T}_k}
\newcommand{\dd}{\,\mathrm{d}}
\newcommand{\cfg}[1]{\texttt{#1}}

\title{Position-update mechanisms in the closed loop:\\
a side-by-side derivation of every implemented method}
\author{Notes accompanying \texttt{closed\_loop.position\_update},
\texttt{closed\_loop.soft\_loss} and\\
\texttt{infl\_ens.inflgame.router.verification}}
\date{August 2026}

\begin{document}
\maketitle

\begin{abstract}
The closed loop alternates \emph{routing} (assigning prompts to agents),
\emph{learning} (LoRA SFT on the assigned prompts) and a
\emph{trait-space position update} (a blended step toward a weighted centroid
of prompts). The influencer-game theory is only ``intact'' when the expected
position drift points along the strategic gradient $\nabla_{x_i}u_i$. This
note writes out, in full, the expected drift of every routing / weighting
combination the code implements --- hard routing (canonical, strategic, with
and without the $(1-G)$ re-weight, expected-pool) and soft top-$k$ routing
(the ``distributed learning'' signal and the ``top-$k$ winners at unit loss''
ablation) --- and shows algebraically which of them are gradient-matched and
why the others are not. A single lemma decides every case: the centroid mass
must be proportional to $B(b)\,G_i(b)\,(1-G_i(b))$ with a proportionality
constant that does not depend on the prompt $b$. Under soft routing that
mass is applied to the \emph{whole} batch, which makes the rule exact for
every $k$ and, by Rao--Blackwell, strictly lower-variance than the
hard-routing update with the same expectation. Every agent --- including a
member of a co-located pair --- steps independently along its own gradient;
merge groups affect only which adapter trains on which prompt, and pairs
persist because identical positions receive identical steps, a prediction
the run measures rather than enforces. The last section reproduces the
numerical cosine and Monte-Carlo checks.
\end{abstract}

\tableofcontents

%======================================================================
\section{Setting and notation}
%======================================================================

\paragraph{Trait space and resource.}
Prompts are embedded into the trait box $[0,1]^L$. A prompt is a point
$b\in[0,1]^L$; the prompt distribution is the resource density $B(b)$,
normalised so that $\sum_b B(b)=1$ over a grid (or $\int B\dd b=1$). A training
round draws a minibatch $b_1,\dots,b_M \stackrel{\text{iid}}{\sim} B$.

\paragraph{Agents and kernel.}
There are $N$ agents with positions $x_1,\dots,x_N\in[0,1]^L$, collected in
$\bx$. Each agent carries a Gaussian kernel with a shared covariance $\Sig$
(isotropic, $\Sig=\sigma^2 I$, unless stated otherwise):
\begin{equation}
  f_i(b) \;=\; \exp\!\Big(-\tfrac12\,(x_i-b)^{\!\top}\Sig^{-1}(x_i-b)\Big).
  \label{eq:kernel}
\end{equation}

\paragraph{Allocation share.}
The (canonical, proportional) allocation of prompt $b$ to agent $i$ is the
softmax over agents,
\begin{equation}
  G_i(\bx,b) \;=\; \frac{f_i(b)}{\sum_{j=1}^N f_j(b)},
  \qquad \sum_{i=1}^N G_i(\bx,b)=1 .
  \label{eq:G}
\end{equation}
In code this is \texttt{allocation\_weights}; every object below is derived
from it.

\paragraph{Utility and its gradient.}
Agent $i$'s utility is the resource it captures,
\begin{equation}
  u_i(\bx) \;=\; \sum_b B(b)\,G_i(\bx,b).
  \label{eq:u}
\end{equation}
Differentiating \eqref{eq:G} with respect to $x_i$ only (the other positions
are held fixed --- this is the \emph{own} gradient that best-response /
gradient-ascent dynamics consume):
\begin{align}
  \nabla_{x_i} f_i(b) &= f_i(b)\,\Sig^{-1}(b-x_i),\\
  \nabla_{x_i} G_i(\bx,b)
    &= \frac{\nabla_{x_i} f_i \cdot \sum_j f_j - f_i\cdot \nabla_{x_i} f_i}{(\sum_j f_j)^2}
     = G_i\,(1-G_i)\;\Sig^{-1}(b-x_i),
  \label{eq:gradG}\\
  \nabla_{x_i} u_i(\bx)
    &= \sum_b B(b)\,G_i(\bx,b)\,\big(1-G_i(\bx,b)\big)\;\Sig^{-1}(b-x_i).
  \label{eq:gradu}
\end{align}
Equation \eqref{eq:gradu} is \texttt{utility\_gradient}. The scalar
coefficient $G_i(1-G_i)$ is the whole story: it is \emph{not} $G_i$, and it is
\emph{not} a per-prompt renormalised version of $G_i$.

\paragraph{Kernel-agnostic form.}
Nothing above uses the Gaussian shape except the last factor. For any
positive differentiable kernel $k(x_i,b)$ with score
$s_i(b):=\nabla_{x_i}\log k(x_i,b)$ the same algebra gives
$\nabla_{x_i}G_i = G_i(1-G_i)\,s_i(b)$ and
$\nabla_{x_i}u_i=\sum_b B(b)G_i(1-G_i)\,s_i(b)$
(\texttt{kernel\_agnostic\_gradient\_step.tex}); for the Gaussian kernel
$s_i(b)=\Sig^{-1}(b-x_i)$, which is what turns a ``score average'' into a
``centroid minus position''. Every statement below is made for the Gaussian
case and holds verbatim with $(b-x_i)\to\Sig\,s_i(b)$ otherwise.

\paragraph{Top-$k$ training support.}
For a prompt $b$ let $\Tk(b)$ be the set of the $k$ agents with the largest
$G_j(\bx,b)$ (ties broken by \texttt{argpartition}). The \emph{renormalised
top-$k$ share} of the update memo is
\begin{equation}
  w_i(b) \;=\; \frac{G_i(b)}{Z_k(b)}\,\ind{i\in\Tk(b)},
  \qquad Z_k(b) := \sum_{j\in\Tk(b)} G_j(b),
  \qquad \sum_i w_i(b)=1 ,
  \label{eq:wtopk}
\end{equation}
(\texttt{top\_k\_allocation\_weights}). In the kept design the top-$k$ set
decides \emph{who trains on what}; it plays no role in the theory-matched
position step (Section~\ref{sec:soft}).

%======================================================================
\section{The common template and the one lemma that decides everything}
\label{sec:template}
%======================================================================

Every method in the code moves agent $i$ by the same rule: pick a
\emph{working set} $S_i$ of prompts and a non-negative \emph{centroid mass}
$m_i(b)$ on it, form the mass-weighted centroid, and blend:
\begin{equation}
  \hat x_i \;=\; \frac{\sum_{b\in S_i} m_i(b)\,b}{\sum_{b\in S_i} m_i(b)},
  \qquad
  x_i \;\leftarrow\; (1-\beta)\,x_i + \beta\,\hat x_i,
  \qquad
  x_i \;\leftarrow\; \operatorname{clip}(x_i,0,1).
  \label{eq:template}
\end{equation}
(\texttt{position\_from\_corpus} / \texttt{soft\_pair\_position\_target} with
\texttt{scores}$=m_i$, then \texttt{apply\_position\_update}; the clip is the
\texttt{clip\_to\_box} option.) The displacement is therefore
\begin{equation}
  \Delta x_i \;=\; \beta\,(\hat x_i - x_i)
  \;=\; \frac{\beta}{\sum_{b\in S_i} m_i(b)}\;
        \underbrace{\sum_{b\in S_i} m_i(b)\,(b-x_i)}_{=:\;d_i} .
  \label{eq:drift}
\end{equation}
The prefactor is a positive scalar, so the \emph{direction} of the step is
the direction of $d_i$. Methods differ only in (i) how $S_i$ is chosen
(sampled or deterministic) and (ii) what $m_i(b)$ is. When $S_i$ is random,
what we can control is the expected mass at each prompt location,
\begin{equation}
  \bar m_i(b) \;:=\; \E\big[m_i(b)\,\ind{b\in S_i}\big]
  \;=\; B(b)\;\Pr[\,b\in S_i\,]\;\cdot\;\text{(weight applied to $b$ once it is in $S_i$)},
  \label{eq:mbar}
\end{equation}
and the expected direction $\bar d_i=\sum_b \bar m_i(b)(b-x_i)$.\footnote{As
in \texttt{kernel\_agnostic\_gradient\_step.tex}, $\hat x_i$ is a ratio of two
Monte-Carlo sums; the law of large numbers in $M$ gives
$d_i/\sum m_i \to \bar d_i/\sum\bar m_i$, and the numerator alone is an
unbiased estimator of $\bar d_i$. ``Matched in expectation'' below always
refers to $\bar d_i$.}

\begin{lemma}[direction matching]
\label{lem:match}
Suppose $\bar m_i(b) = c_i\,B(b)\,G_i(b)\,(1-G_i(b))$ for some constant
$c_i>0$ that does not depend on $b$. Then
\[
  \bar d_i \;=\; c_i\,\Sig\,\nabla_{x_i}u_i(\bx),
\]
so for isotropic $\Sig=\sigma^2 I$ the expected step is \emph{exactly
parallel} to $\nabla_{x_i}u_i$, and for general $\Sig$ it is the
$\Sig$-preconditioned gradient (still an ascent direction, since
$\Sig\succ0$).
\end{lemma}
\begin{proof}
Substitute into $\bar d_i=\sum_b \bar m_i(b)(b-x_i)$ and compare with
\eqref{eq:gradu}: $\sum_b c_i B G_i(1-G_i)(b-x_i)
= c_i\,\Sig\sum_b B G_i(1-G_i)\Sig^{-1}(b-x_i)= c_i\Sig\nabla_{x_i}u_i$.
\end{proof}

\begin{remark}[what breaks it]
\label{rem:break}
Write $\bar m_i(b) = B(b)G_i(1-G_i)\,\rho_i(b)$. The lemma needs
$\rho_i(b)\equiv c_i$. Two distinct ways to violate it appear in the
implemented methods:
\begin{enumerate}
\item \textbf{Missing factor.} $\rho_i(b) = 1/(1-G_i(b))$, i.e.\ the mass is
  $B\,G_i$ instead of $B\,G_i(1-G_i)$. The difference
  $\sum_b B\,G_i^{\,2}\,(b-x_i)$ is dominated by the agent's own stronghold
  (where $G_i\approx1$), so the step is pulled toward the centre of mass of
  what the agent already owns rather than toward the contested boundary where
  the gradient lives.
\item \textbf{Per-prompt normaliser.} $\rho_i(b) = 1/Z(b)$ for some
  $Z(b)$ that varies with $b$ (a softmax over agents evaluated at $b$, or a
  top-$k$ renormaliser). Prompts where the normaliser is small are
  over-weighted relative to the gradient. Renormalising \emph{per agent}
  (dividing the whole sum by $\sum_b m_i(b)$, as \eqref{eq:template} does) is
  harmless because that constant does not depend on $b$.
\end{enumerate}
A third potential defect --- restricting $S_i$ to a subset of prompts on
which $G_i(1-G_i)$ is not identically zero --- is a truncation bias. It
never appears in the kept design, because the position step always uses
either the whole batch or the whole pool (Section~\ref{sec:soft}).
\end{remark}

Everything that follows is an application of Lemma~\ref{lem:match} and
Remark~\ref{rem:break} to each method's $\bar m_i$.

%======================================================================
\section{Hard routing: one prompt per (sampled) winning agent}
\label{sec:hard}
%======================================================================

Under \cfg{routing\_mode: hard} each prompt $b_m$ is assigned to exactly one
agent $a_m$ drawn from a categorical distribution $p_\cdot(b_m)$ over agents
(\cfg{policy: proportional}; \cfg{argmax} replaces the draw by
$\arg\max_i p_i$). The working set is $S_i=\{m:\,a_m=i\}$ and a per-prompt
centroid weight $\omega_i(b)$ is applied inside it. Hence
\begin{equation}
  \Pr[b\in S_i]=p_i(b),\qquad
  \bar m_i(b) \;=\; B(b)\,p_i(b)\,\omega_i(b).
  \label{eq:hardmass}
\end{equation}
The sampling step contributes the factor $p_i(b)$ ``for free''; the centroid
weight $\omega_i$ must supply whatever is missing.

\subsection{H1: canonical routing, uniform centroid (\cfg{position\_update: naive})}
$p_i=G_i$, $\omega_i\equiv1$:
\begin{equation}
  \bar m_i(b) = B(b)\,G_i(b),\qquad
  \bar d_i = \sum_b B\,G_i\,(b-x_i)\quad\text{(\emph{not} matched: missing factor).}
\end{equation}
This is the historical default (uniform mean of routed embeddings). It is the
``canonical naive'' column of the verification report. Configs recorded with
it are now pinned to \cfg{position\_update: naive}.

\subsection{H2: canonical routing, $(1-G)$ centroid (\cfg{position\_update: theory\_matched}, the default)}
$p_i=G_i$, $\omega_i(b)=1-G_i(b)$:
\begin{equation}
  \bar m_i(b) = B(b)\,G_i(b)\,(1-G_i(b)),\qquad
  \bar d_i = \Sig\,\nabla_{x_i}u_i\quad\text{(matched exactly, $c_i=1$).}
\end{equation}
The SFT loss runs at unit weight; only the centroid is re-weighted. This is
what the deprecated \cfg{loss\_reweight: position\_only} used to select and
is also precisely what the no-SFT simulator
\texttt{simulate\_position\_only\_loop} realises
(\texttt{weights\_i = 1 - G\_batch[i, mine\_idx]}). ``Position-only'' and
``theory-matched hard routing'' are therefore the same expected dynamics; the
former simply skips the LoRA.

\subsection{H3: full re-weight (\cfg{loss\_reweight: one\_minus\_G})}
Same $\bar m_i$ as H2 (matched), but the per-example SFT loss weight is also
$s_i(b)=1-G_i(b)$. The learning signal then has effective sample size
\[
  \mathrm{ESS}_i=\frac{\big(\sum_{m\in S_i}(1-G_i(b_m))\big)^2}{\sum_{m\in S_i}(1-G_i(b_m))^2}
  \;\le\; |S_i|,
\]
which is small once the agent owns a stronghold ($1-G_i\approx0$ on most of
its prompts). H2 keeps the full ESS; that is the reason the loss side and the
centroid side are decoupled knobs.

\subsection{H4: strategic routing (\cfg{routing\_weight: G\_times\_1mG})}
$p_i(b)=\dfrac{G_i(1-G_i)}{S(b)}$ with $S(b):=\sum_j G_j(b)(1-G_j(b))$,
$\omega_i\equiv1$:
\begin{equation}
  \bar m_i(b) = B(b)\,\frac{G_i(1-G_i)}{S(b)},\qquad
  \rho_i(b)=\frac{1}{S(b)}\quad\text{(approximately matched: per-prompt normaliser).}
\end{equation}
The normaliser $S(b)$ is a genuine function of $b$ (it is large in contested
regions and tiny inside strongholds), so by Remark~\ref{rem:break}(2) the
direction is only approximately the gradient --- the ``strategic'' column of
the report reads $0.98$--$0.999$ rather than $1.000$. Under
\cfg{theory\_matched} the centroid stays \emph{uniform} here: routing already
carries the whole coefficient, and applying $(1-G_i)$ a second time would
give $\bar m_i\propto B\,G_i(1-G_i)^2/S(b)$, which is why the validator
rejects \cfg{G\_times\_1mG} together with any $(1-G)$ re-weight.

\subsection{H5: expected-pool centroid (\cfg{centroid\_mode: expected\_pool})}
Instead of the sampled working set, use the whole pool deterministically with
mass $m_i(b_m)=G_i(1-G_i)$:
\begin{equation}
  \hat x_i=\frac{\sum_m G_i(b_m)(1-G_i(b_m))\,b_m}{\sum_m G_i(b_m)(1-G_i(b_m))},
  \label{eq:pool}
\end{equation}
which is the $M\to\infty$ limit of H2 (\texttt{expected\_pool\_centroid}).
Matched by construction; it has no meaningful ``naive'' variant, so the
validator rejects \cfg{expected\_pool} with \cfg{position\_update: naive}.
It is available under both routing modes.

%======================================================================
\section{Soft (dense, top-$k$) routing}
\label{sec:soft}
%======================================================================

Under \cfg{routing\_mode: soft} nothing is sampled. Each prompt is assigned
\emph{deterministically} to all agents in $\Tk(b)$ for \emph{training}:
\begin{equation}
  S_i^{\mathrm{train}}=\{m:\,i\in\Tk(b_m)\} .
  \label{eq:softS}
\end{equation}
The crucial difference from Section~\ref{sec:hard}: \textbf{there is no
sampling factor $p_i(b)$.} Whatever coefficient the gradient needs must be
put into the centroid weight $\omega_i$ explicitly. The second crucial
observation is that the position step need not use $S_i^{\mathrm{train}}$
at all: it is a weighted mean of coordinates that are already projected
for the whole batch, so using every prompt costs nothing.

\subsection{S1: renormalised-share centroid (\cfg{position\_update: naive}) --- the memo's rule}
The update memo (steps 1--4) uses the renormalised top-$k$ share
\eqref{eq:wtopk} as the centroid weight on the training set,
$\omega_i(b)=w_i(b)$, $S_i=S_i^{\mathrm{train}}$:
\begin{equation}
  \bar m_i(b)=B(b)\,\frac{G_i(b)}{Z_k(b)}\,\ind{i\in\Tk(b)},
  \qquad
  \rho_i(b)=\frac{\ind{i\in\Tk(b)}}{(1-G_i(b))\,Z_k(b)} .
  \label{eq:S1}
\end{equation}
Every defect of Remark~\ref{rem:break} occurs at once:
\begin{itemize}
\item the $(1-G_i)$ factor is missing --- even in the fully dense case
  $k=N$ (where $Z_N\equiv1$) the mass is $B\,G_i$, i.e.\ \emph{exactly the
  canonical-naive direction H1};
\item for $k<N$ the top-$k$ renormaliser $Z_k(b)$ additionally rescales each
  prompt by a $b$-dependent amount, and the indicator truncates the sum.
\end{itemize}
Writing the dense case out,
\[
  \bar d_i^{\,\mathrm{S1}} = \sum_b B\,G_i\,(b-x_i)
  = \underbrace{\sum_b B\,G_i(1-G_i)(b-x_i)}_{\Sig\nabla_{x_i}u_i}
   + \underbrace{\sum_b B\,G_i^{\,2}\,(b-x_i)}_{\text{stronghold pull}} ,
\]
and the second term can dominate and even reverse the direction: in the
verification configuration with one agent per resource mode, the naive
cosine is $-1.000$ while the matched one is $+1.000$ (Section~\ref{sec:num}).
This is the ``distributed learning'' arm as originally implemented, and it is
why the answer to ``do the position-only and the distributed-learning steps
match in gradient direction?'' is \emph{no}. It is kept only as the
explicit ablation arm.

\subsection{S2: dense $G(1-G)$ mass over the whole batch (\cfg{position\_update: theory\_matched}, the default)}
\label{sec:S2}
Put the full gradient coefficient into the centroid weight, apply it to
\emph{every} prompt of the batch, and do \emph{not} renormalise per prompt:
\begin{equation}
  S_i=\{1,\dots,M\},\qquad
  m_i(b_m) = G_i(b_m)\,(1-G_i(b_m)),\qquad
  \bar m_i(b) = B(b)\,G_i(b)\,(1-G_i(b)) .
  \label{eq:S2}
\end{equation}
This is \texttt{matched\_centroid\_mass} applied to the batch allocation
matrix. Lemma~\ref{lem:match} applies with $c_i=1$ for \emph{any}
\cfg{soft\_top\_k}:
\begin{proposition}
\label{prop:S2}
Under S2 the expected drift is $\bar d_i^{\,\mathrm{S2}}=\Sig\,\nabla_{x_i}u_i$
exactly, independently of $k$. It coincides with the position-only /
theory-matched hard dynamics H2 in expectation, and with the expected-pool
centroid \eqref{eq:pool} in the limit $M\to\infty$.
\end{proposition}
\begin{proof}
The indicator of \eqref{eq:S1} is gone and $\rho_i\equiv1$. Since
$\Pr[b\in S_i]=1$, \eqref{eq:mbar} gives $\bar m_i(b)=B(b)G_i(1-G_i)$.
\end{proof}

\begin{proposition}[Rao--Blackwell: S2 has lower variance than H2]
\label{prop:RB}
Fix the batch $b_1,\dots,b_M$. Let $d_i^{\mathrm{H2}}=\sum_m \ind{a_m=i}\,(1-G_i(b_m))\,(b_m-x_i)$
be the numerator of the hard rule (with the categorical draw $a_m\sim G_\cdot(b_m)$)
and $d_i^{\mathrm{S2}}=\sum_m G_i(b_m)(1-G_i(b_m))\,(b_m-x_i)$ that of the dense rule. Then
\[
  \E\big[d_i^{\mathrm{H2}}\,\big|\,b_{1:M}\big]=d_i^{\mathrm{S2}},
  \qquad
  \Var\big(d_i^{\mathrm{S2}}\big)\le\Var\big(d_i^{\mathrm{H2}}\big),
\]
with strict inequality whenever some $G_i(b_m)\in(0,1)$.
\end{proposition}
\begin{proof}
$\E[\ind{a_m=i}\mid b_m]=G_i(b_m)$, so conditioning the hard numerator on
the batch gives the dense one term by term. The variance statement is the
law of total variance,
\[
  \Var\big(d_i^{\mathrm{H2}}\big)
  =\Var\big(\E[d_i^{\mathrm{H2}}\mid b_{1:M}]\big)+\E\big[\Var(d_i^{\mathrm{H2}}\mid b_{1:M})\big]
  =\Var\big(d_i^{\mathrm{S2}}\big)+\E\Big[\sum_m G_i(1-G_i)\,(1-G_i)^2\,\|b_m-x_i\|^2\Big],
\]
and the second term is positive unless every $G_i(b_m)\in\{0,1\}$.
\end{proof}
So ``just use the hard routing's update'' would be correct but wasteful:
the dense rule uses all $M$ prompts for every agent instead of the
$\approx G_i M$ it happened to win, at the same expectation. The
Monte-Carlo table in Section~\ref{sec:num} shows the spread roughly halved.

\begin{remark}[why the mass must not be renormalised per prompt]
Replacing \eqref{eq:S2} by $B\,G_i(1-G_i)/Z'(b)$ with
$Z'(b)=\sum_{j}G_j(1-G_j)$ (or any other $b$-dependent normaliser) would
give $\rho_i(b)=1/Z'(b)$ and reintroduce Remark~\ref{rem:break}(2). The
per-\emph{agent} normalisation $\sum_{m}m_i(b_m)$ in \eqref{eq:template} is a
single positive number and is harmless.
\end{remark}
\begin{remark}[the training gate is irrelevant to the position step]
Because $S_i=\{1,\dots,M\}$, an agent that trained on no prompt this round
(possible for small $k$) still takes its gradient step; and an agent that
trained on every prompt ($k\ge N$) takes exactly the same step it would with
$k=1$. The top-$k$ gate bounds SFT cost; it does not enter the theory.
\end{remark}

\subsection{S3: co-located merge groups (pairs) --- training only}
\label{sec:S3}
With \cfg{sft\_merge\_groups} a group $p$ of clones shares one LoRA. Define
the group share
\begin{equation}
  G_p(\bx,b) := \sum_{i\in p} G_i(\bx,b)
  \qquad(\texttt{group\_allocation\_weights}).
\end{equation}
If every group has the same size $s$ and its members are co-located at
$x_p$, then $f_i=f_p$ for $i\in p$ and
\[
  G_p(\bx,b)=\frac{s\,f_p(b)}{\sum_q s\,f_q(b)}=\frac{f_p(b)}{\sum_q f_q(b)},
\]
i.e.\ exactly the allocation of one agent at $x_p$ in the $P$-player game.
That identity is what justifies taking the top-$k$ \emph{training} gate over
groups rather than clones.

\paragraph{Positions are not grouped.}
The group share is used for \emph{nothing else}. Every clone takes its own
S2 step in the full $N$-player game, with its own $G_i$ --- its twin
appearing as one competitor among the others:
\begin{equation}
  m_i(b_m)=G_i(b_m)\big(1-G_i(b_m)\big)\ \text{ over the whole batch,
  independently for each } i .
  \label{eq:S3}
\end{equation}
No shared step is ever written; the trainer's routing position merely tracks
the mean of its members. The reason this is the right design is that the
theory already predicts co-location:

\begin{proposition}[co-located clones stay co-located, unforced]
\label{prop:pairs}
If $x_i=x_j$ then $G_i(\bx,b)=G_j(\bx,b)$ for every $b$, hence
$m_i\equiv m_j$, hence $\hat x_i=\hat x_j$ and the two clones receive
\emph{identical} updates in \eqref{eq:template}. Co-location is therefore
invariant under the independent dynamics: it is a fixed point of the pair's
relative coordinate, not a constraint imposed by the update.
\end{proposition}
\begin{proof}
$G_i$ depends on $i$ only through $f_i$, which depends only on $x_i$; equal
positions give equal kernels, hence equal shares, masses, targets and blends.
\end{proof}

\noindent The practical consequence is that co-location becomes a
\emph{measurement} rather than an assumption. Each round logs the within-pair
distance $\|x_i-x_j\|$ in \texttt{agent\_geometry}; if the theory's stability
claim were to fail (a partner drifting off under numerical noise, an
asymmetric blend cap, a $\sigma$ above the threshold), the run would show it
instead of hiding it behind a shared write. The verification report confirms
the mechanism numerically: the twin drifts differ by exactly $0$
(Section~\ref{sec:num}, configuration D).

The historical pair rule survives only as the ablation arm
(\cfg{position\_update: naive}): one step per group, computed from the
renormalised group share on the pair's top-$k$ prompts and written to both
members. It inherits every defect of S1.

%======================================================================
\section{The loss side: what the agents actually train on}
\label{sec:loss}
%======================================================================

The learning signal is independent of the centroid mass. With per-example
weights $s_i(b)$ the SFT objective is the weight-normalised mean of the
response-token cross-entropy,
\begin{equation}
  \mathcal{L}_i(\theta_i)=\frac{1}{|S_i^{\mathrm{train}}|}\sum_{m\in S_i^{\mathrm{train}}} s_i(b_m)\,\ell(\theta_i;b_m),
  \qquad \ell=\text{mean token NLL of the formatted (prompt, response)} ,
\end{equation}
so $s_i$ acts as a per-example learning-rate multiplier
(\texttt{weighted\_causal\_lm\_loss}); $s_i\equiv1$ takes the stock unweighted
trainer path. The implemented choices are:
\begin{center}
\begin{tabular}{@{}lll@{}}
\toprule
routing & config & $s_i(b)$ on the training set \\
\midrule
hard & \cfg{loss\_reweight: null} (default) & $1$ \\
hard & \cfg{loss\_reweight: one\_minus\_G} & $1-G_i(b)$ \\
soft & \cfg{soft\_loss: weighted} (default) & $w_i(b)=G_i(b)/Z_k(b)$ \quad(``distributed'' signal) \\
soft & \cfg{soft\_loss: unit} & $1$ \quad(``top-$k$ winners'' ablation) \\
\bottomrule
\end{tabular}
\end{center}
With \cfg{sft\_merge\_groups} the group replaces the clone as the unit of
this table ($G_p$, $w_p$, one adapter per group); the position step is
unaffected (Section~\ref{sec:S3}).

\paragraph{The top-$k$ winners ablation.}
Hard routing gives each prompt to \emph{one} winner at unit weight. The
ablation \cfg{routing\_mode: soft}, \cfg{soft\_loss: unit} gives each prompt
to its $k$ highest-share agents, \emph{each at unit weight} --- the natural
top-$k$ generalisation of one-prompt-per-winning-agent, with $k=1$ recovering
a deterministic (argmax rather than sampled) winner. Compared with the
weighted soft signal it removes the share weighting from the learning signal
while leaving the assignment untouched; compared with hard routing it
replaces one sampled winner by $k$ deterministic ones. Because the centroid
mass is chosen separately (and, under \cfg{theory\_matched}, is the same
whole-batch mass for every $k$), the ablation grid is
\[
  \{\text{weighted},\ \text{unit}\}\times\{\text{theory\_matched},\ \text{naive}\}
  \ (\times\ k),
\]
and every arm logs \texttt{soft\_loss}, \texttt{position\_update},
\texttt{soft\_top\_k}, the batch, and the applied centroid weights
(\texttt{agent\_position\_weights}, aligned with \texttt{batch\_prompts}) per
round.

\paragraph{Data volume.}
Under unit-weight top-$k$ each specialist sees up to $k\times$ the tokens
that the de-duplicated pooled replay sees (the replay keeps one copy of each
source prompt per round). The comparator is data-matched at the level of
source prompts, not of gradient steps; capability comparisons should be read
with that asymmetry in mind.

%======================================================================
\section{Summary table}
%======================================================================

$\bar m_i(b)$ is the expected centroid mass at prompt $b$
(Section~\ref{sec:template}); ``matched'' means Lemma~\ref{lem:match} holds
with a $b$-independent constant.

\begin{center}
\footnotesize
\setlength{\tabcolsep}{4pt}
\begin{tabular}{@{}>{\raggedright\arraybackslash}p{2.6cm}>{\raggedright\arraybackslash}p{4.3cm}>{\raggedright\arraybackslash}p{3.5cm}>{\raggedright\arraybackslash}p{1.5cm}>{\raggedright\arraybackslash}p{3.3cm}@{}}
\toprule
method & config keys & $\bar m_i(b)/B(b)$ & loss weight $s_i$ & matched? \\
\midrule
H1 canonical naive & \cfg{routing\_mode: hard}, \cfg{routing\_weight: G}, \cfg{position\_update: naive} & $G_i$ on routed prompts & $1$ & no (missing $1-G_i$) \\[3pt]
H2 theory-matched (default) & \cfg{hard}, \cfg{G}, \cfg{position\_update: theory\_matched} & $G_i(1-G_i)$ on routed prompts & $1$ & \textbf{exact} \\[3pt]
H3 full re-weight & \cfg{hard}, \cfg{G}, \cfg{loss\_reweight: one\_minus\_G} & $G_i(1-G_i)$ on routed prompts & $1-G_i$ & \textbf{exact} (low ESS) \\[3pt]
H4 strategic & \cfg{hard}, \cfg{routing\_weight: G\_times\_1mG} & $G_i(1-G_i)/S(b)$ & $1$ & $\approx$ (per-prompt normaliser) \\[3pt]
H5 expected pool & \cfg{centroid\_mode: expected\_pool} (either routing mode) & $G_i(1-G_i)$ over the pool & any & \textbf{exact} \\[3pt]
S1 soft naive & \cfg{routing\_mode: soft}, \cfg{position\_update: naive} & $G_i\,\ind{i\in\Tk}/Z_k(b)$ & $w_i$ or $1$ & no (all defects) \\[3pt]
S2 soft matched (default) & \cfg{soft}, \cfg{position\_update: theory\_matched} & $G_i(1-G_i)$ over the whole batch & $w_i$ or $1$ & \textbf{exact}, any $k$; lower variance than H2 \\[3pt]
S3 soft pairs & \cfg{soft} $+$ \cfg{sft\_merge\_groups} & S2 per clone; $G_p$ used only for the training gate & $w_p$ or $1$ & \textbf{exact} per clone; pairs persist unforced (Prop.~\ref{prop:pairs}) \\
\bottomrule
\end{tabular}
\end{center}

The position-only simulator is H2 without SFT. Hence: \emph{the position-only
update and the soft ``distributed learning'' update agree in expected
direction if and only if the soft arm uses S2; the original soft arm (S1)
does not, and S2 is additionally the lower-variance estimator of the two.}

%======================================================================
\section{Numerical verification}
\label{sec:num}
%======================================================================

\texttt{run\_reweighted\_drift\_report} evaluates \eqref{eq:gradu} and each
$\bar d_i$ on a $25\times25$ grid over a bimodal resource landscape (modes at
$(0.25,0.25)$ and $(0.75,0.75)$, width $0.18$; kernel $\sigma=0.20$) and
prints the cosine similarity $\cos(\bar d_i,\nabla_{x_i}u_i)$. \emph{dense
match} is S2 (it does not depend on $k$); \emph{top$k$ naive} is S1 at the
given $k$. Configuration D is the co-located-pairs check of
Section~\ref{sec:S3}. The values below are from the run recorded on
2026-08-26.

\begin{center}
\footnotesize
\setlength{\tabcolsep}{3.5pt}
\resizebox{\textwidth}{!}{%
\begin{tabular}{@{}llrrrrrr@{}}
\toprule
configuration & agent & strategic (H4) & canon+$(1{-}G)$ (H2) & canon naive (H1) & dense match (S2) & top3 naive (S1) & top2 naive (S1) \\
\midrule
A: near-symmetric Nash & 0 & 0.9994 & 1.0000 & 0.4947 & \textbf{1.0000} & 0.4947 & 0.0887 \\
\quad $x=(.5,.4),(.5,.55),(.45,.5)$ & 1 & 0.9935 & 1.0000 & 0.9090 & \textbf{1.0000} & 0.9090 & 0.8718 \\
 & 2 & 0.9821 & 1.0000 & 0.5587 & \textbf{1.0000} & 0.5587 & $-0.6336$ \\
\midrule
B: one agent per mode & 0 & 1.0000 & 1.0000 & $-1.0000$ & \textbf{1.0000} & $-1.0000$ & $-1.0000$ \\
\quad $+$ contested centre & 1 & 1.0000 & 1.0000 & $-1.0000$ & \textbf{1.0000} & $-1.0000$ & $-1.0000$ \\
\quad $x=(.28,.28),(.72,.72),(.5,.5)$ & 2$^{\dagger}$ & 0.0025 & $-0.0051$ & 0.0023 & $-0.0072$ & 0.0022 & 0.0046 \\
\midrule
C: two on one mode, one alone & 0--2 & 1.0000 & 1.0000 & 1.0000 & 1.0000 & 1.0000 & 1.0000 \\
\bottomrule
\end{tabular}}

\medskip
\begin{tabular}{@{}llrrr@{}}
\toprule
\multicolumn{5}{@{}l}{D: two co-located clone pairs, four clones, cosines against each
  \emph{clone's own} $\nabla_{x_i}u_i$ in the 4-player game} \\
\midrule
pair & clone & naive pair rule ($k=2$) & S3 per-clone match & $\|d_{\text{twin}}-d_i\|$ \\
\midrule
0 at $(0.30,0.35)$ & 0 & 0.9535 & \textbf{1.0000} & $0.00\mathrm{e}{+}00$ \\
                   & 1 & 0.9535 & \textbf{1.0000} & $0.00\mathrm{e}{+}00$ \\
1 at $(0.70,0.60)$ & 2 & 0.9749 & \textbf{1.0000} & $0.00\mathrm{e}{+}00$ \\
                   & 3 & 0.9749 & \textbf{1.0000} & $0.00\mathrm{e}{+}00$ \\
\bottomrule
\end{tabular}
\end{center}

\noindent Reading the table:
\begin{itemize}
\item \textbf{top3 naive $=$ canon naive}, column for column. This is the
  identity $Z_N\equiv1$ in \eqref{eq:S1}: fully dense soft routing with the
  share centroid is the canonical-naive rule.
\item \textbf{dense match $=1.0000$ everywhere the gradient is non-zero},
  as Proposition~\ref{prop:S2} states; it equals the canon$+(1-G)$ column,
  i.e.\ the position-only dynamics, and does so for every $k$.
\item In B the naive direction is \emph{anti}-parallel to the gradient for
  the two mode-holders: their stronghold pull \eqref{eq:S1} points back
  into the mode while the gradient points toward the contested centre.
\item $^{\dagger}$Agent 2 in B sits at a stationary point
  ($\|\nabla u_2\|\approx10^{-16}$); its cosines are numerical noise, not
  evidence for or against any rule.
\item In C all drifts are collinear (every rule pushes along the diagonal
  joining the modes), so the direction test is uninformative there.
\item In D the last column is the whole point of Proposition~\ref{prop:pairs}:
  the two members of a pair compute their steps completely independently
  and the results differ by \emph{exactly} zero, so the pair holds together
  with nothing enforcing it. Each clone's own drift is parallel to its own
  gradient; the naive pair rule (one shared step from the renormalised
  group share) is close but not parallel here --- and it is the arm that
  cannot detect a pair coming apart, since it writes one position to both
  members by construction.
\end{itemize}

\paragraph{Finite-batch variance (Proposition~\ref{prop:RB}).}
The same report draws 200 batches of 4096 prompts, realises each rule on the
sampled batch (hard rules with an actual categorical routing draw), and
reports the cosine of the \emph{mean} drift with the gradient together with
the standard deviation of the drift norm across batches:

\begin{center}
\footnotesize
\begin{tabular}{@{}llrrr@{}}
\toprule
configuration & rule & cos(mean drift, $\nabla u_i$) & $\|$mean drift$\|$ & std of $\|$drift$\|$ \\
\midrule
A, agent 0 & H2 canonical$+(1{-}G)$ & 0.9998 & 0.0507 & 0.0052 \\
           & S2 dense matched        & 1.0000 & 0.0508 & \textbf{0.0026} \\
A, agent 1 & H2 & 0.9999 & 0.0491 & 0.0050 \\
           & S2 & 1.0000 & 0.0494 & \textbf{0.0024} \\
A, agent 2 & H2 & 0.9999 & 0.0281 & 0.0048 \\
           & S2 & 1.0000 & 0.0282 & \textbf{0.0025} \\
\midrule
B, agent 0 & H2 & 1.0000 & 0.0916 & 0.0042 \\
           & S2 & 1.0000 & 0.0919 & \textbf{0.0023} \\
B, agent 1 & H2 & 1.0000 & 0.0920 & 0.0043 \\
           & S2 & 1.0000 & 0.0919 & \textbf{0.0023} \\
\midrule
C, agent 2 & H2 & 1.0000 & 0.2812 & 0.0062 \\
           & S2 & 1.0000 & 0.2814 & \textbf{0.0034} \\
\bottomrule
\end{tabular}
\end{center}

\noindent The means agree to three decimals (same expectation) and the
spread of the dense rule is roughly half that of the hard rule, exactly as
the law-of-total-variance decomposition predicts. The naive rules are
omitted here; their mean drifts point elsewhere entirely (the cosine table
above).

The same facts are asserted by \texttt{tests/test\_topk\_matched.py}
(dense matched cosine $>0.9999$; naive strictly lower at $k=N$ and $k=2$;
co-located clones take bitwise identical independent steps; dense spread
$<$ hard spread) and end-to-end in stubbed closed-loop rounds (hard default
applies $(1-G)$; soft default applies the whole-batch $G(1-G)$ mass
\emph{per clone} and skips SFT's own position update; \cfg{soft\_loss: unit}
passes no loss weights; partners end each round at the same position with a
logged within-pair distance of $0$; \cfg{expected\_pool} works under soft
routing).

%======================================================================
\section{Configuration cheat-sheet}
%======================================================================

\begin{verbatim}
closed_loop:
  routing_mode: hard | soft          # sampled single winner | deterministic top-k
  routing_weight: G | G_times_1mG    # canonical | strategic   (hard only)
  soft_top_k: 2                      # soft only; who trains on what
  soft_loss: weighted | unit         # soft only; share-weighted | top-k winners
  loss_reweight: null | one_minus_G  # hard only, loss side; 'position_only'
                                     #   is a deprecated alias for the default
  position_update: theory_matched | naive   # centroid mass (default: matched)
  centroid_mode: batch | expected_pool      # whole batch | whole pool
\end{verbatim}

\noindent\textbf{theory\_matched} applies, per mode, the mass that makes
Lemma~\ref{lem:match} hold: $(1-G_i)$ on the routed prompts under hard
canonical routing, uniform under strategic routing (the sampling already
carries $G_i(1-G_i)$), and the dense $G_i(1-G_i)$ over the whole batch under
soft routing, independent of \cfg{soft\_top\_k} and of
\cfg{sft\_merge\_groups} --- every agent steps on its own.
\textbf{naive} restores the historical centroid (uniform mean under hard
routing, renormalised top-$k$ share under soft routing, one shared step per
merge group) as an explicit ablation arm.

\end{document}
